% Reviewed translation: PDF pages 82--87 / printed pages 68--73.

\MathBookSourcePage{82}
\begin{Proof}
证明的论证与 Theorem~\ref{thm:abstract-penalty-error} 非常相似. 置
\[
w^\varepsilon=u^\varepsilon-u-\sum_{n=1}^N\varepsilon^nu_n,
\qquad
\rho^\varepsilon=\lambda^\varepsilon-\lambda-
\sum_{n=1}^N\varepsilon^n\lambda_n.
\]
利用 
\eqref{eq:abstract-q-first},
\eqref{eq:abstract-q-second},
\eqref{eq:abstract-regularized-q-first},
\eqref{eq:abstract-regularized-q-second} 和
\eqref{eq:abstract-asymptotic-recursion}, 得
\begin{equation}\label{eq:abstract-asymptotic-error-system}
\left\{
\begin{aligned}
a(w^\varepsilon,v)+b(v,\rho^\varepsilon)&=0
&&\forall v\in X,\\
-\varepsilon c(\rho^\varepsilon,\mu)+b(w^\varepsilon,\mu)
&=\varepsilon^{N+1}c(\lambda_N,\mu)
&&\forall\mu\in M.
\end{aligned}
\right.
\end{equation}
第一个方程 
\eqref{eq:abstract-asymptotic-error-system} 与
\eqref{eq:abstract-inf-sup} 给出
\begin{equation}\label{eq:abstract-asymptotic-rho-bound}
\|\rho^\varepsilon\|_M
\le(\|a\|/\beta)\|w^\varepsilon\|_X.
\end{equation}
接着在 
\eqref{eq:abstract-asymptotic-error-system} 中取
$v=w^\varepsilon$ 和 $\mu=\rho^\varepsilon$, 得
\[
\begin{aligned}
a(w^\varepsilon,w^\varepsilon)
&=-\varepsilon c(\rho^\varepsilon,\rho^\varepsilon)
-\varepsilon^{N+1}c(\lambda_N,\rho^\varepsilon)\\
&\le-\varepsilon^{N+1}c(\lambda_N,\rho^\varepsilon),
\end{aligned}
\]
所以
\[
a(w^\varepsilon,w^\varepsilon)
\le\varepsilon^{N+1}(\|a\|\|c\|/\beta)
\|\lambda_N\|_M\|w^\varepsilon\|_X.
\]
另一方面, 
\eqref{eq:abstract-asymptotic-error-system} 中第二个方程给出
\[
Bw^\varepsilon=\varepsilon C(\rho^\varepsilon+\varepsilon^N\lambda_N),
\]
所以
\[
\begin{aligned}
\langle C^{-1}Bw^\varepsilon,Bw^\varepsilon\rangle
&=\varepsilon^2c(\rho^\varepsilon+\varepsilon^N\lambda_N,
\rho^\varepsilon+\varepsilon^N\lambda_N)\\
&\le2\varepsilon^2\|c\|
(\|\rho^\varepsilon\|_M^2+\varepsilon^{2N}\|\lambda_N\|_M^2),
\end{aligned}
\]
并由 
\eqref{eq:abstract-asymptotic-rho-bound},
\[
\langle C^{-1}Bw^\varepsilon,Bw^\varepsilon\rangle
\le2\varepsilon^2\|c\|
\{(\|a\|/\beta)^2\|w^\varepsilon\|_X^2
+\varepsilon^{2N}\|\lambda_N\|_M^2\}.
\]
因此利用 
\eqref{eq:abstract-combined-ellipticity}, 得到形如
\[
\alpha\|w^\varepsilon\|_X^2
\le C_1\varepsilon^2\|w^\varepsilon\|_X^2
+C_2\varepsilon^{N+1}\|\lambda_N\|_M\|w^\varepsilon\|_X
+C_3\varepsilon^{2N+2}\|\lambda_N\|_M^2
\]
的不等式. 当 $\varepsilon\le\varepsilon_0$ 足够小时, 这蕴含
\begin{equation}\label{eq:abstract-asymptotic-w-bound}
\|w^\varepsilon\|_X
\le C_4\varepsilon^{N+1}\|\lambda_N\|_M.
\end{equation}
现在由 
\eqref{eq:abstract-asymptotic-recursion} 和
Corollary~\ref{cor:abstract-v-elliptic-well-posedness},
\begin{equation}\label{eq:abstract-asymptotic-lambda-n-bound}
\|\lambda_N\|_M\le C_5\|\lambda\|_M
\le C_6(\|l\|_{X'}+\|\chi\|_{M'}),
\end{equation}
其中常数 $C_5$ 和 $C_6$ 具有 $C^N$ 的形式.
所需不等式 
\eqref{eq:abstract-penalty-asymptotic-bound} 因而由
\eqref{eq:abstract-asymptotic-rho-bound},
\eqref{eq:abstract-asymptotic-w-bound} 和
\eqref{eq:abstract-asymptotic-lambda-n-bound} 得到.
\end{Proof}

\setcounter{subsection}{3}
\subsection{梯度型迭代方法}\label{subsec:abstract-gradient-iterations}

本小节的目的是导出求解问题 $(Q)$ 的便捷迭代方法.
为此, 我们提出一种基于优化表述
\eqref{eq:abstract-constrained-minimization} 和
\eqref{eq:abstract-dual-minimization} 的对偶方法.

\MathBookSourcePage{83}
暂且假设双线性形式 $b(\cdot,\cdot)$ 满足 inf-sup 条件
\eqref{eq:abstract-inf-sup}, 且双线性形式 $a(\cdot,\cdot)$ 对称并在
\eqref{eq:abstract-x-ellipticity} 的意义下 $X$-椭圆.
由 Corollary~\ref{cor:abstract-v-elliptic-well-posedness}, 问题 $(Q)$ 适定;
Theorem~\ref{thm:abstract-saddle-equivalence} 断言其唯一解 $(u,\lambda)$
由 
\eqref{eq:abstract-lagrangian} 定义的 Lagrange 泛函
$\mathcal L(v,\mu)$ 的唯一鞍点刻画. 此外, $u$ 是约束优化问题
\eqref{eq:abstract-constrained-minimization} 的唯一解, 而 $\lambda$ 是无约束优化问题
\eqref{eq:abstract-dual-minimization} 的唯一解. 显然最后这个问题最容易求解,
而下降法应当特别适合求解它. 这里集中讨论一个一般下降算法,
并把它应用于两类梯度法: 简单梯度法(也称 Uzawa Algorithm)和共轭梯度法.

基于上述考虑, 先推广 Theorem~\ref{thm:abstract-saddle-equivalence},
说明 $(u,\lambda)$ 可由一族参数化 Lagrange 泛函的唯一鞍点刻画.
按照上一小节, 引入连续双线性形式
$c(\cdot,\cdot):M\times M\to\mathbb R$, 它在
\eqref{eq:abstract-c-ellipticity} 的意义下 $M$-椭圆且对称.
同样, 用 
\eqref{eq:abstract-operator-c} 定义相应算子
$C\in\mathcal L(M;M')$. 对所有 $r\ge0$, 定义增广能量泛函
$J_r:X\to\mathbb R$ 和增广 Lagrange 泛函
$\mathcal L_r:X\times M\to\mathbb R$:
\begin{equation}\label{eq:abstract-augmented-energy}
J_r(v)=J(v)+(r/2)\langle C^{-1}(Bv-\chi),Bv-\chi\rangle,
\end{equation}
\begin{equation}\label{eq:abstract-augmented-lagrangian}
\mathcal L_r(v,\mu)=J_r(v)+b(v,\mu)-\langle\chi,\mu\rangle.
\end{equation}
由于
\[
J_r(v)=J(v)\qquad\forall v\in V(\chi),
\]
$u$ 仍由下式刻画:
\begin{equation}\label{eq:abstract-augmented-constrained-minimum}
J_r(u)=\inf_{v\in V(\chi)}J_r(v).
\end{equation}
另一方面, 有如下结果.

\setcounter{statement}{4}
\begin{Theorem}\label{thm:abstract-augmented-saddle}
假设 
\eqref{eq:abstract-inf-sup},
\eqref{eq:abstract-v-ellipticity} 和
\eqref{eq:abstract-c-ellipticity} 成立, 并再假设双线性形式
$a(\cdot,\cdot)$ 和 $c(\cdot,\cdot)$ 对称, 且
\begin{equation}\label{eq:abstract-augmented-semipositive}
a(v,v)+r\langle C^{-1}Bv,Bv\rangle\ge0
\qquad\forall v\in X.
\end{equation}
则问题 $(Q)$ 的解 $(u,\lambda)$ 是增广 Lagrange 泛函
$\mathcal L_r$ 的唯一鞍点, 或等价地由下式刻画:
\begin{equation}\label{eq:abstract-augmented-minmax}
\operatorname{Min}_{v\in X}
\left(\sup_{\mu\in M}\mathcal L_r(v,\mu)\right)
=\mathcal L_r(u,\lambda)
=\operatorname{Max}_{\mu\in M}
\left(\inf_{v\in X}\mathcal L_r(v,\mu)\right).
\end{equation}
\end{Theorem}
\begin{Proof}
考虑 $a(\cdot,\cdot)$ 与 $c(\cdot,\cdot)$ 的对称性, 得
\[
D_v\mathcal L_r(u,\lambda)\mathbin{:}v
=a(u,v)-\langle l,v\rangle
+r\langle C^{-1}(Bu-\chi),Bv\rangle+b(v,\lambda),
\]
\[
D_{vv}^2\mathcal L_r(u,\lambda)\mathbin{:}(v,v)
=a(v,v)+r\langle C^{-1}Bv,Bv\rangle.
\]
于是证明沿用 Theorem~\ref{thm:abstract-saddle-equivalence} 和
Corollary~\ref{cor:abstract-saddle-minmax} 的论证.
\end{Proof}

\MathBookSourcePage{84}
现在转向 
\eqref{eq:abstract-augmented-constrained-minimum} 的对偶问题
\[
\operatorname{Max}_{\mu\in M}\inf_{v\in X}\mathcal L_r(v,\mu).
\]
定义双线性形式
\[
a_r(u,v)=a(u,v)+r\langle C^{-1}Bu,Bv\rangle.
\]
假设它是 $X$-椭圆的, 即存在 $\alpha_r>0$, 使
\begin{equation}\label{eq:abstract-ar-ellipticity}
a_r(v,v)\ge\alpha_r\|v\|_X^2
\qquad\forall v\in X.
\end{equation}
此时显然优化问题
\[
\mathcal L_r(u_r(\mu),\mu)=\inf_{v\in X}\mathcal L_r(v,\mu)
\]
有唯一解 $u_r(\mu)\in X$. 事实上, 根据上述证明, 这个问题等价于求解
\begin{equation}\label{eq:abstract-ur-state-equation}
a_r(u_r(\mu),v)=\langle l,v\rangle
+r b(v,C^{-1}\chi)-b(v,\mu)
\qquad\forall v\in X.
\end{equation}
仍有 $u_r(\lambda)=u$, 因而
\begin{equation}\label{eq:abstract-augmented-dual-maximum}
\mathcal L_r(u_r(\lambda),\lambda)
=\operatorname{Max}_{\mu\in M}\mathcal L_r(u_r(\mu),\mu).
\end{equation}
仿照 Subsection~\ref{subsec:abstract-saddle-point}, 容易得到
\[
\mathcal L_r(u_r(\mu),\mu)
=-(1/2)a_r(u_r(\mu),u_r(\mu))-\langle\chi,\mu\rangle
+(r/2)\langle C^{-1}\chi,\chi\rangle.
\]
因此置
\begin{equation}\label{eq:abstract-augmented-dual-functional}
K_r(\mu)=(1/2)a_r(u_r(\mu),u_r(\mu))+\langle\chi,\mu\rangle,
\end{equation}
便有
\[
\mathcal L_r(u_r(\mu),\mu)
=-K_r(\mu)+(r/2)\langle C^{-1}\chi,\chi\rangle.
\]
由这个等式和 
\eqref{eq:abstract-augmented-dual-maximum}, 得 $\lambda$ 满足
\begin{equation}\label{eq:abstract-augmented-dual-minimum}
K_r(\lambda)=\operatorname{Min}_{\mu\in M}K_r(\mu).
\end{equation}
下面的 Corollary 概括了这个结果.

\setcounter{statement}{2}
\begin{Corollary}\label{cor:abstract-augmented-dual-characterization}
在 Theorem~\ref{thm:abstract-augmented-saddle} 的假设和椭圆性条件
\eqref{eq:abstract-ar-ellipticity} 下, $\lambda\in M$ 由无约束优化问题
\eqref{eq:abstract-augmented-dual-minimum} 的唯一解刻画.
\end{Corollary}

\setcounter{statement}{4}
\begin{Remark}\label{rem:abstract-ur-affine}
由 
\eqref{eq:abstract-ur-state-equation} 定义的映射
$\mu\mapsto u_r(\mu)$ 是从 $M$ 到 $X$ 的仿射映射.
\end{Remark}

为了用下降法求解 
\eqref{eq:abstract-augmented-dual-minimum},
必须计算二次泛函 $K_r$ 关于 $\mu$ 的前两个导数.
首先由 
\eqref{eq:abstract-ur-state-equation} 看出 $u_r$ 的导数

\MathBookSourcePage{85}
\[
Du_r=Du_r(\mu)\in\mathcal L(M;X)
\]
与 $\mu$ 无关. 更确切地, 
\eqref{eq:abstract-ur-state-equation} 给出
\[
Du_r=-A_r^{-1}B',
\]
其中 $A_r\in\mathcal L(X;X')$ 表示与双线性形式 $a_r(\cdot,\cdot)$
相联系的算子:
\[
\langle A_ru,v\rangle=a_r(u,v)
\qquad\forall u,v\in X.
\]
因此对 
\eqref{eq:abstract-augmented-dual-functional} 求导, 得
\[
DK_r(\mu)\mathbin{:}v
=\langle A_rDu_r\mathbin{:}v,u_r(\mu)\rangle+\langle\chi,v\rangle
=\langle\chi-Bu_r(\mu),v\rangle.
\]
所以
\begin{equation}\label{eq:abstract-kr-first-derivative}
DK_r(\mu)=\chi-Bu_r(\mu)\in M',
\end{equation}
以及
\begin{equation}\label{eq:abstract-kr-second-derivative}
D^2K_r(\mu)=D^2K_r=-BDu_r=BA_r^{-1}B'
\in\mathcal L(M;M').
\end{equation}

接着给 Hilbert 空间 $M$ 配备一个内积. 由于双线性形式
$c(\cdot,\cdot)$ 对称且 $M$-椭圆, 可以选它作为内积.
相对于 $c$, $K_r$ 在 $M$ 中点 $\mu$ 处的梯度 $g_r(\mu)$ 定义为
\[
c(g_r(\mu),v)=DK_r(\mu)\mathbin{:}v.
\]
换言之,
\begin{equation}\label{eq:abstract-kr-gradient}
g_r(\mu)=C^{-1}(\chi-Bu_r(\mu)).
\end{equation}

现在设 $(\omega^m)$ 是 $M$ 中任意元素序列, 考虑如下下降法:
从初值 $\lambda^0\in M$ 出发, 按下式构造序列 $(\lambda^m)\subset M$:
\begin{equation}\label{eq:abstract-general-descent-lambda}
\lambda^{m+1}=\lambda^m-\rho^m\omega^m,
\qquad m\ge0,
\end{equation}
其中参数 $\rho^m$ 的选择使得
\[
K_r(\lambda^m-\rho^m\omega^m)
=\inf_{\rho\in\mathbb R}K_r(\lambda^m-\rho\omega^m).
\]
最优 $\rho^m$ 为
\begin{equation}\label{eq:abstract-optimal-step}
\rho^m=\frac{DK_r(\lambda^m)\mathbin{:}\omega^m}
{D^2K_r\mathbin{:}(\omega^m,\omega^m)},
\end{equation}
而 $u^m$ 当然由 
\eqref{eq:abstract-ur-state-equation} 与 $\lambda^m$ 联系:
\begin{equation}\label{eq:abstract-um-from-lambda}
u^m=A_r^{-1}(l+rB'C^{-1}\chi-B'\lambda^m)=u_r(\lambda^m).
\end{equation}
注意, 这种迭代方法把 $u^m$ 与 $\lambda^m$ 的计算解耦.
序列 $(\omega^m)$ 目前尚未指定. 后面将看到, 适当选择 $(\omega^m)$
会导出很有吸引力的梯度法. 从实际观点看, 算法
\eqref{eq:abstract-general-descent-lambda}--\eqref{eq:abstract-um-from-lambda}
可以组织得更高效.

\MathBookSourcePage{86}
为简单起见, 置
\[
g^m=g_r(\lambda^m),
\qquad
z^m=A_r^{-1}B'\omega^m.
\]
由 
\eqref{eq:abstract-kr-second-derivative},
\eqref{eq:abstract-kr-gradient} 和
\eqref{eq:abstract-optimal-step}, 容易得到
\[
\rho^m=c(g^m,\omega^m)/b(z^m,\omega^m).
\]
此外, 将 
\eqref{eq:abstract-um-from-lambda} 的两个相邻值相减并利用
\eqref{eq:abstract-general-descent-lambda}, 得
\[
u^{m+1}-u^m=\rho^mz^m.
\]
因此一般下降算法具有如下步骤:

1) 给定初值 $\lambda^0\in M$, 由下式计算 $u^0\in X$:
\[
A_ru^0=l+B'(rC^{-1}\chi-\lambda^0).
\]

2) 对 $m\ge0$, 给定 $\omega^m\in M$, 并已知
$(u^m,\lambda^m)\in X\times M$, 由下式确定
$(z^m,g^m)\in X\times M$, $\rho^m\in\mathbb R$ 以及
$(u^{m+1},\lambda^{m+1})\in X\times M$:
\begin{equation}\label{eq:abstract-general-descent-system}
\left\{
\begin{aligned}
Cg^m&=\chi-Bu^m,\\
A_rz^m&=B'\omega^m,\\
\rho^m&=c(g^m,\omega^m)/b(z^m,\omega^m),\\
\lambda^{m+1}&=\lambda^m-\rho^m\omega^m,\\
u^{m+1}&=u^m+\rho^mz^m.
\end{aligned}
\right.
\end{equation}
每次迭代必须求解两个线性问题: 一个使用算子 $C$, 另一个使用算子 $A_r$.
实际中算子 $C$ 容易求逆. 例如对 Stokes 问题, 通常选
$L^2$ 内积作为 $c(\cdot,\cdot)$, 使 $C$ 为恒等算子.
所以主要计算工作集中在 $z^m$ 上.

这个下降法有如下收敛结果.

\setcounter{statement}{5}
\begin{Theorem}\label{thm:abstract-general-descent-convergence}
假设形式 $b(\cdot,\cdot)$ 满足 inf-sup 条件
\eqref{eq:abstract-inf-sup}, 且 $a_r(\cdot,\cdot)$ 与
$c(\cdot,\cdot)$ 都对称, 并分别满足椭圆性条件
\eqref{eq:abstract-ar-ellipticity} 与
\eqref{eq:abstract-c-ellipticity}. 则对每个初值 $\lambda^0\in M$
和每个序列 $(\omega^m)\subset M$, 迭代格式
\eqref{eq:abstract-general-descent-system} 定义 $X\times M$ 中唯一序列
$(u^m,\lambda^m)$. 此外, 如果序列 $(g^m)$ 与 $(\omega^m)$ 对所有
$m\ge0$ 满足
\begin{equation}\label{eq:abstract-descent-angle-condition}
c(g^m,\omega^m)\ge\alpha\|g^m\|_M\|\omega^m\|_M,
\qquad \alpha>0\text{ independent of }m,
\end{equation}
则下降法收敛:
\[
\lim_{m\to\infty}
\{\|u^m-u\|_X+\|\lambda^m-\lambda\|_M\}=0.
\]
\end{Theorem}
\begin{Proof}
显然, 由于 $a_r(\cdot,\cdot)$ 和 $c(\cdot,\cdot)$ 的椭圆性,
\eqref{eq:abstract-general-descent-system} 中前两个方程有唯一解.
此外, 
\eqref{eq:abstract-kr-second-derivative} 给出

\MathBookSourcePage{87}
\[
D^2K_r\mathbin{:}(\mu,\mu)
=\langle A_r^{-1}B'\mu,B'\mu\rangle.
\]
所以 inf-sup 条件 
\eqref{eq:abstract-inf-sup} 蕴含
\begin{equation}\label{eq:abstract-kr-hessian-bounds}
\delta_r\|\mu\|_M^2
\le D^2K_r\mathbin{:}(\mu,\mu)
\le\tau_r\|\mu\|_M^2
\qquad\forall\mu\in M,quad\delta_r,\tau_r>0.
\end{equation}
因此 
\eqref{eq:abstract-general-descent-system} 定义唯一的 $\rho^m$
(当然假设 $\omega^m\ne0$).

现在讨论算法的收敛性. 首先由构造有
\[
DK_r(\lambda^{m+1})\mathbin{:}\omega^m=0.
\]
由于 $K_r$ 为二次泛函, 这蕴含
\[
K_r(\lambda^m)-K_r(\lambda^{m+1})
=(1/2)D^2K_r\mathbin{:}
(\lambda^m-\lambda^{m+1},\lambda^m-\lambda^{m+1}).
\]
于是由 
\eqref{eq:abstract-kr-hessian-bounds},
\[
K_r(\lambda^m)-K_r(\lambda^{m+1})
\ge(1/2)\delta_r\|\lambda^m-\lambda^{m+1}\|_M^2.
\]
但由构造, 序列 $K_r(\lambda^m)$ 单调递减, 且以下界
$K_r(\lambda)$ 有界. 所以它收敛, 从而
\[
\lim_{m\to\infty}\|\lambda^m-\lambda^{m+1}\|_M=0.
\]
由 
\eqref{eq:abstract-general-descent-lambda}, 这意味着
\[
\lim_{m\to\infty}\rho^m\|\omega^m\|_M=0.
\]
而 
\eqref{eq:abstract-kr-hessian-bounds} 与假设
\eqref{eq:abstract-descent-angle-condition} 蕴含 $(\rho^m)$ 有下界
\[
\rho^m\ge(\alpha/\tau_r)
(\|g^m\|_M/\|\omega^m\|_M).
\]
因此
\[
\rho^m\|\omega^m\|_M
\ge(\alpha/\tau_r)\|g^m\|_M,
\]
从而
\[
\lim_{m\to\infty}g^m=0,
\]
即
\begin{equation}\label{eq:abstract-dual-gradient-convergence}
\lim_{m\to\infty}DK_r(\lambda^m)=0
\quad\text{in }M'.
\end{equation}
此外, 因为 $DK_r(\lambda)=0$, 可写成
\[
DK_r(\lambda^m)=DK_r(\lambda^m)-DK_r(\lambda)
=D^2K_r(\lambda^m-\lambda),
\]
\[
D^2K_r\mathbin{:}(\lambda^m-\lambda,\lambda^m-\lambda)
=DK_r(\lambda^m)\mathbin{:}(\lambda^m-\lambda).
\]
由 
\eqref{eq:abstract-kr-hessian-bounds}, 得
\[
\delta_r\|\lambda^m-\lambda\|_M
\le\|DK_r(\lambda^m)\|_{M'}.
\]
% Proof continues on PDF page 88 / printed page 74.
